\documentclass[11pt,a4paper]{article}
\input{preamble}
\title{Principles of Statistics}
\begin{document}
\maketitle
\thispagestyle{empty}
\section*{Principles of Statistics}

\subsection*{0. INTRODUCTION}

Let $(\Omega,\mathcal{F},\mathbb{P})$ be a probability space, and
$X:\Omega\to\mathbb{R}$ a random variable with distribution function
\[
\begin{aligned}
F(t)
&=
\mathbb{P}\bigl(\{\omega:X(\omega)\leq t\}\bigr) \\
&=
\begin{cases}
\displaystyle \int_{-\infty}^{t} f(x)\,dx,
& \text{if $X$ is continuous with pdf $f$}, \\[1ex]
\displaystyle \sum_{x\leq t} f(x),
& \text{if $X$ is discrete with pmf $f$}.
\end{cases}
\end{aligned}
\]

We denote $P=\mathbb{P}\circ X^{-1}$, the law or ``distribution'' of $X$.

We observe $X_1,\ldots,X_n$, which are iid copies of $X$, where $n$ is
the \underline{sample size}.

\paragraph{Def 1.}
A \underline{statistical model} is a collection
$\{f_\theta:\theta\in\Theta\}$ of pdfs or pmfs, or
$\{P_\theta:\theta\in\Theta\}$ of laws, indexed by a
\underline{parameter space} $\Theta$.

\paragraph{Ex.}
\begin{itemize}
    \item $N(\theta,1)$, \qquad $\Theta=\mathbb{R}$.
    \item $N(\mu,\sigma^2)$, \qquad
    $\theta=(\mu,\sigma^2)$, \quad
    $\Theta=\mathbb{R}\times(0,\infty)$.
    \item $\operatorname{Poi}(\theta)$, \qquad
    $\Theta=[0,\infty)$.
\end{itemize}

\paragraph{Def 2.}
We say the model is \underline{correctly specified} when
$\exists\,\theta_0\in\Theta$ such that $P=P_{\theta_0}$.

We will refer to the ``true parameter'' as $\theta_0$.

\subsection*{Goals of Statistics}

\begin{enumerate}
    \item \underline{Estimation}:
    We define an \underline{estimator}
    \[
    \widehat{\theta}_n
    =
    \widehat{\theta}_n(X_1,\ldots,X_n),
    \]
    which we hope is close to $\theta_0$ when
    $X_i\overset{\mathrm{iid}}{\sim}P_{\theta_0}$,
    for all $\theta_0\in\Theta$, when $n$ is large.

    \item \underline{Testing}:
    We have a null hypothesis
    \[
    H_0:\theta\in\Theta_0\subseteq\Theta,
    \]
    and an alternative hypothesis
    \[
    H_1:\theta\in\Theta_1\subseteq\Theta.
    \]
    We define a \underline{test},
    \[
    \psi_n:=\psi_n(X_1,\ldots,X_n),
    \]
    such that $\psi_n=0$ if $H_0$ is true, and
    $\psi_n=1$ if $H_1$ is true, \underline{w.h.p.}

    \item \underline{Uncertainty Quantification}
    \begin{itemize}
        \item Define a \underline{confidence set}
        $C_n(X_1,\ldots,X_n;\alpha)$, such that
        \[
        P_\theta(\theta\in C_n)
        \overset{(\geq)}{=}
        1-\alpha
        \qquad \forall\,\theta\in\Theta,
        \]
        where $\alpha$ is the \underline{significance level}.

        \item Define a \underline{Bayesian credible set},
        a region of high posterior probability.
    \end{itemize}
\end{enumerate}

\section*{I. LIKELIHOOD PRINCIPLE}

\paragraph{Def 3.}
The \underline{likelihood function}
\[
L_n(\theta)=\prod_{i=1}^{n}f_\theta(X_i).
\]

\noindent\underline{Note:}
This is a random function, as it depends implicitly on $X_1,\ldots,X_n$.

The \underline{log-likelihood function}
\[
\ell_n(\theta)=\sum_{i=1}^{n}\log f_\theta(X_i),
\]
we write
\[
\overline{\ell}_n(\theta)=\frac{1}{n}\ell_n(\theta)
\]
for a \underline{normalised log-likelihood}.

\paragraph{Def 4.}
A \underline{maximum likelihood estimator}
$\widehat{\theta}_{\mathrm{MLE}}$ is any estimator with
\[
L_n\bigl(\widehat{\theta}_{\mathrm{MLE}}\bigr)
=
\max_{\theta\in\Theta}L_n(\theta).
\]

\paragraph{Ex.}
\[
X_i\overset{\mathrm{iid}}{\sim}\operatorname{Poi}(\theta).
\]
\[
L_n(\theta)
=
\prod_{i=1}^{n}e^{-\theta}\frac{\theta^{X_i}}{X_i!}.
\]
\[
\ell_n(\theta)
=
-n\theta+(\log\theta)\sum X_i+\mathrm{const.}
\]
\[
\ell_n'(\theta)
=
-n+\frac{\sum X_i}{\theta}
\overset{\mathrm{FOC}}{=}0
\]
\[
\Longrightarrow\quad
\text{suggests }
\widehat{\theta}_{\mathrm{MLE}}
=
\frac{\displaystyle\sum_{i=1}^{n}X_i}{n}
:=
\overline{X}_n
\qquad
\text{``sample mean''}.
\]

\noindent\underline{SOC:}
\[
\ell_n''(\theta)
=
-\frac{1}{\theta^2}\sum X_i
<0
\qquad
\text{if }\sum X_i>0.
\]

If $\sum X_i=0$,
\[
\ell_n(\theta)=-n\theta
\quad\Longrightarrow\quad
\widehat{\theta}_{\mathrm{MLE}}=0.
\]

In every case
\[
\widehat{\theta}_{\mathrm{MLE}}=\overline{X}_n.
\]

\paragraph{Ex.}
\[
X_i\overset{\mathrm{iid}}{\sim}N(\mu,\sigma^2),
\qquad
\theta=(\mu,\sigma^2).
\]
\[
\ell_n(\theta)
=
-\frac{n}{2}\log\pi
-\frac{n}{2}\log(\sigma^2)
-\frac{1}{2\sigma^2}
\sum_{i=1}^{n}(X_i-\mu)^2.
\]

\noindent\underline{FOC:}
\[
\frac{\partial\ell_n}{\partial\mu}
=
-\frac{1}{\sigma^2}
\sum_{i=1}^{n}(X_i-\mu)
=0,
\]
\[
\frac{\partial\ell_n}{\partial\sigma^2}
=
-\frac{n}{2\sigma^2}
+
\frac{1}{2(\sigma^2)^2}
\sum_{i=1}^{n}(X_i-\mu)^2
=0.
\]

We can check that the Hessian of $\ell_n$ is negative definite
$\forall\,\theta\in\mathbb{R}\times(0,\infty)$.

Hence,
\[
\widehat{\theta}_{\mathrm{MLE}}
=
\begin{pmatrix}
\overline{X}_n \\[6pt]
\underbrace{
\displaystyle\frac{1}{n}
\sum_{i=1}^{n}(X_i-\overline{X}_n)^2
}_{\text{``Sample variance''}}
\end{pmatrix}.
\]

\subsection*{Non-examinable}

The likelihood principle poses that if two observations
\[
X=(X_1,\ldots,X_n)
\qquad\text{and}\qquad
X^*=(X_1^*,\ldots,X_n^*)
\]
give rise to likelihoods $L_n,L_n^*$ respectively, such that
\[
L_n(\theta)=L_n^*(\theta)\cdot M
\qquad \forall\,\theta\in\Theta,
\]
and some constant $M$, then they should produce the same
``inference'' on the true parameter.

\begin{itemize}
    \item The MLE, Bayes estimators, \underline{credible sets}
    satisfy the likelihood principle.
    \item \underline{Confidence sets}, tests generally do not.
\end{itemize}

\noindent\underline{Birnbaum (1962)}
\[
\text{Conditionality}+\text{Sufficiency}
\quad\Longrightarrow\quad
\text{Likelihood principle}.
\]

\begin{center}
\textbf{end of lecture 1}
\end{center}

\section*{II. Information Geometry}

The \underline{score function} is
\[
S_n(\theta)
=
\frac{\partial}{\partial\theta}\ell_n(\theta),
\qquad
\Theta\subseteq\mathbb{R}^d.
\]

In many cases,
\[
S_n(\widehat{\theta}_{\mathrm{MLE}})=0,
\]
where, by convention,
\[
S_n(\theta')
=
\left.
\frac{\partial}{\partial\theta}\ell_n(\theta)
\right|_{\theta=\theta'}.
\]

Let $X$ be a r.v.\ taking values in $\mathcal{X}$, and
$g:\mathcal{X}\to\mathbb{R}$.
Take a model $\{f(x,\theta):\theta\in\Theta\}$.
We will write
\[
\mathbb{E}_{\theta}g(X)
=
\int g(x)f(x,\theta)\,dx,
\qquad
X\sim f(\cdot,\theta).
\]

\medskip
\noindent\textbf{Obs 1.}

Suppose the model $\{f(\cdot,\theta):\theta\in\Theta\}$ is such that
\[
\mathbb{E}_{\theta_0}\bigl|\log f(X,\theta)\bigr|
<
\infty
\qquad
\forall\theta\in\Theta.
\]
Define
\[
\ell(\theta)
=
\mathbb{E}_{\theta_0}\log f(X,\theta),
\qquad
X\sim P_{\theta_0}.
\]

The normalised log-likelihood
\[
\overline{\ell}_n(\theta)
=
\frac{1}{n}\sum_{i=1}^{n}\log f(X_i,\theta)
\]
can be thought of as a sample version of $\ell(\theta)$.

Suppose the model is correctly specified, such that $X$ has pdf
$f(\cdot,\theta_0)$ for some $\theta_0\in\Theta$.

\medskip
\noindent\textbf{Fact.}
Then $\ell(\theta)$ is maximised at $\theta_0$.
This justifies using the MLE to approximate $\theta_0$ when $n$ is large.

Indeed,
\[
\begin{aligned}
\ell(\theta)-\ell(\theta_0)
&=
\mathbb{E}_{\theta_0}
\left(
\log\frac{f(X,\theta)}{f(X,\theta_0)}
\right)
\\
&\leq
\log\mathbb{E}_{\theta_0}
\left(
\frac{f(X,\theta)}{f(X,\theta_0)}
\right)
\\
&=
\log\int
\frac{f(x,\theta)}{f(x,\theta_0)}
f(x,\theta_0)\,dx
\\
&=
\log 1=0.
\end{aligned}
\]
By Jensen's inequality, as $\log$ is concave, with equality if
\[
\frac{f(X,\theta)}{f(X,\theta_0)}
\]
is a.s.\ constant.

\[
\therefore\quad
\ell(\theta)\leq\ell(\theta_0).
\]

If
\[
f(\cdot,\theta)=f(\cdot,\theta')
\quad\Longrightarrow\quad
\theta=\theta',
\]
we say the parameter is \underline{identifiable}.

In that case, if $\theta\neq\theta_0$, Jensen's inequality is strict,
so $\theta_0$ is actually the unique maximizer of $\ell$.

\medskip
\noindent\textbf{Rmk.}
The difference
\[
\ell(\theta_0)-\ell(\theta)
=
\int
\log\frac{f(x,\theta_0)}{f(x,\theta)}
f(x,\theta_0)\,dx
\]
is called the \underline{Kullback--Leibler divergence}
or \underline{relative entropy} between $P_{\theta_0}$ and $P_\theta$.

A notion of distance which is \underline{not} symmetric.

\medskip
\noindent\textbf{Obs 2.}

\[
\begin{aligned}
0
&=
\frac{\partial}{\partial\theta}
\underbrace{\int f(x,\theta)\,dx}_{=1\ \forall\theta}
=
\int\frac{\partial}{\partial\theta}f(x,\theta)\,dx
\\
&=
\int
\frac{\partial\log f(x,\theta)}{\partial\theta}
f(x,\theta)\,dx
\\
&=
\mathbb{E}_{\theta}
\left[
\frac{\partial\log f(X,\theta)}{\partial\theta}
\right].
\end{aligned}
\]

[There are some conditions for this.
(Leibniz rule). See P \& M.]

``$S_n(\theta)$ is centred when $X\sim P_\theta$.''

\medskip
\noindent\textbf{Obs 3.}

\medskip
\noindent\textbf{Def 5.}
Let $\Theta\subseteq\mathbb{R}^d$.
For any $\theta\in\operatorname{int}(\Theta)$,
the \underline{Fisher Information matrix} is
\[
I(\theta)
=
\mathbb{E}_{\theta}
\left[
\frac{\partial}{\partial\theta}\log f(X,\theta)
\cdot
\left(
\frac{\partial}{\partial\theta}\log f(X,\theta)
\right)^T
\right],
\qquad
X\sim P_\theta.
\]

\medskip
\noindent\textbf{Rmk.}
This is the covariance matrix of
\[
S_n(\theta)
=
\frac{\partial}{\partial\theta}\log f(X,\theta)
\]
when $X\sim P_\theta$.

\medskip
\noindent\textbf{Prop 1.}
For all $\theta\in\operatorname{int}(\Theta)$, in a ``regular'' model,
\[
I(\theta)
=
-\mathbb{E}_{\theta}
\left[
\frac{\partial^2}{\partial\theta\,\partial\theta^T}
\log f(X,\theta)
\right]
\]
\[
\left(
I_{ij}(\theta)
=
-\mathbb{E}_{\theta}
\left[
\frac{\partial^2}{\partial\theta_i\,\partial\theta_j}
\log f(X,\theta)
\right]
\right).
\]

\medskip
\noindent\textbf{Proof.}
\[
0
=
\frac{\partial^2}{\partial\theta\,\partial\theta^T}
\int f(x,\theta)\,dx
=
\int
\frac{\partial^2}{\partial\theta\,\partial\theta^T}
f(x,\theta)\,dx.
\]

[When this equality holds, we say the model is ``regular''.]

Now, we have
\[
\begin{aligned}
\frac{\partial^2}{\partial\theta\,\partial\theta^T}
\log f(x,\theta)
&=
\frac{\partial}{\partial\theta^T}
\left[
\frac{1}{f(x,\theta)}
\frac{\partial}{\partial\theta}f(x,\theta)
\right]
\\
&=
-\frac{1}{f(x,\theta)^2}
\frac{\partial}{\partial\theta}f(x,\theta)
\left(
\frac{\partial}{\partial\theta}f(x,\theta)
\right)^T
\\
&\qquad
+
\frac{1}{f(x,\theta)}
\frac{\partial^2}{\partial\theta\,\partial\theta^T}
f(x,\theta).
\end{aligned}
\]

Take $-\mathbb{E}_{\theta}[\cdot]$ of both sides.
\[
\begin{aligned}
&-\mathbb{E}_{\theta}
\left[
\frac{\partial^2}{\partial\theta\,\partial\theta^T}
\log f(X,\theta)
\right]
\\
&\qquad=
\mathbb{E}_{\theta}
\left[
\frac{1}{f(X,\theta)^2}
\frac{\partial}{\partial\theta}f(X,\theta)
\left(
\frac{\partial}{\partial\theta}f(X,\theta)
\right)^T
\right]
+0
\\
&\qquad=
\mathbb{E}_{\theta}
\left[
\frac{\partial}{\partial\theta}\log f(X,\theta)
\left(
\frac{\partial}{\partial\theta}\log f(X,\theta)
\right)^T
\right]
=
I(\theta).
\end{aligned}
\]
\hfill$\square$

\medskip
\noindent\textbf{Tensorisation.}

Suppose the observation $X$ consists of $n$ random variables
$X_1,\ldots,X_n$. Then, the F.I.\ matrix is
\[
I_n(\theta)
=
\mathbb{E}_{\theta}
\left[
\frac{\partial}{\partial\theta}
\log f(X_1,\ldots,X_n,\theta)
\left(
\frac{\partial}{\partial\theta}
\log f(X_1,\ldots,X_n,\theta)
\right)^T
\right].
\]

If $X_1,\ldots,X_n$ are iid such that
\[
f(X_1,\ldots,X_n,\theta)
=
\prod_{i=1}^{n}f(X_i,\theta),
\]
where $\{f(x,\theta):\theta\in\Theta\}$ is a model for $X_1$,
then the F.I.\ is
\[
\begin{aligned}
I_n(\theta)
&=
\mathbb{E}_{\theta}
\left[
\sum_{i,j}
\underbrace{
\frac{\partial}{\partial\theta}\log f(X_i,\theta)
}_{\substack{\text{mean }0\text{ under}\\P_\theta}}
\left(
\frac{\partial}{\partial\theta}\log f(X_j,\theta)
\right)^T
\right]
\\
&=
\mathbb{E}_{\theta}
\left[
\sum_{i=1}^{n}
\frac{\partial}{\partial\theta}\log f(X_i,\theta)
\left(
\frac{\partial}{\partial\theta}\log f(X_i,\theta)
\right)^T
\right]
\\
&=
nI(\theta),
\end{aligned}
\]
where $I(\theta)$ is the Fisher Information for a single observation.

\begin{center}
\textbf{End of Lecture 2}
\end{center}

\subsection*{Recap of Information Geometry}

\paragraph{Observation 1.}
\[
\operatorname{KL}(P_{\theta_0},P_\theta)
=
\mathbb{E}_{\theta_0}
\left[
\log\frac{f(X,\theta_0)}{f(X,\theta)}
\right]
\geq 0,
\]
minimised at $\theta=\theta_0$.

\paragraph{Observation 2.}
\[
\mathbb{E}_{\theta}
\left[
\frac{\partial}{\partial\theta}\log f(X,\theta)
\right]
=0,
\qquad
\forall\theta\in\operatorname{int}(\Theta).
\]

\paragraph{Observation 3.}
\[
I(\theta)
:=
\operatorname{Var}_{\theta}
\left[
\frac{\partial}{\partial\theta}\log f(X,\theta)
\right]
=
-\mathbb{E}_{\theta}
\left[
\frac{\partial^2}{\partial\theta\,\partial\theta^{T}}
\log f(X,\theta)
\right],
\qquad
\forall\theta\in\operatorname{int}(\Theta).
\]

\paragraph{Why geometry?}
The F.I., $I(\theta)$, defines a metric tensor on the manifold
\[
\{f(\cdot,\theta):\theta\in\Theta\}.
\]
This is known as the \textbf{Fisher--Rao metric}.
KL divergence is approximated locally by this metric.

\begin{itemize}
    \item \textbf{Reparametrisation invariance.}
    Take the $\operatorname{Bernoulli}(\theta)$ model, where $\theta$
    is the probability of success. Reparametrise it in terms of
    ``log-odds'':
    \[
    \eta=\log\left(\frac{\theta}{1-\theta}\right)\in\mathbb{R}.
    \]
    The manifold
    \[
    \{f(\cdot,\eta):\eta\in\mathbb{R}\},
    \]
    equipped with the metric tensor $I(\eta)$, preserves geodesics,
    path lengths, etc.

    \item \textbf{Sufficient statistic invariance.}
\end{itemize}

\paragraph{Chentsov's theorems.}
These properties characterise the Fisher--Rao metric.

\subsection*{Proposition 2 (Cram\'er--Rao Lower Bound)}

Let
\[
\{f(\cdot,\theta):\theta\in\Theta\subseteq\mathbb{R}\}
\]
be a ``regular'' statistical model.
Let $\tilde{\theta}(X_1,\ldots,X_n)$ be an unbiased estimator based on
i.i.d.\ samples $X_1,\ldots,X_n$:
\[
\mathbb{E}_{\theta}\tilde{\theta}=\theta,
\qquad
\forall\theta\in\Theta.
\]
Then
\[
\operatorname{Var}_{\theta}(\tilde{\theta})
\geq
\frac{1}{nI(\theta)},
\qquad
\forall\theta\in\operatorname{int}(\Theta).
\]

\paragraph{Proof.}
Recall the Cauchy--Schwarz inequality for random variables $Y,Z$:
\[
\bigl(\mathbb{E}[YZ]\bigr)^2
\leq
\mathbb{E}[Y^2]\mathbb{E}[Z^2].
\]
Let
\[
Z=\tilde{\theta}-\theta,
\qquad
Y=\frac{\partial}{\partial\theta}
\log f(X_1,\ldots,X_n;\theta).
\]
Then
\[
\mathbb{E}[Z^2]
=
\operatorname{Var}_{\theta}(\tilde{\theta}),
\qquad
\text{as $\tilde{\theta}$ is unbiased},
\]
and
\[
\mathbb{E}[Y^2]
=
nI(\theta),
\qquad
\text{as the F.I.\ tensorises}.
\]
Moreover,
\[
\begin{aligned}
\mathbb{E}[YZ]
&=
\mathbb{E}_{\theta}
\left[
\tilde{\theta}
\frac{\partial}{\partial\theta}
\log f(X_1,\ldots,X_n;\theta)
\right]
\\
&\quad
-\theta
\underbrace{
\mathbb{E}_{\theta}
\left[
\frac{\partial}{\partial\theta}
\log f(X_1,\ldots,X_n;\theta)
\right]
}_{=0}
\\
&=
\int_{\mathcal{X}^n}
\tilde{\theta}(x_1,\ldots,x_n)
\frac{\partial f(x_1,\ldots,x_n;\theta)}{\partial\theta}
\frac{f(x_1,\ldots,x_n;\theta)}
     {f(x_1,\ldots,x_n;\theta)}
\,dx_1\cdots dx_n
\\
&=
\frac{\partial}{\partial\theta}
\underbrace{
\int_{\mathcal{X}^n}
\tilde{\theta}\,
f(x_1,\ldots,x_n;\theta)
\,dx_1\cdots dx_n
}_{=\mathbb{E}_{\theta}\tilde{\theta}=\theta}
\qquad
\text{(by regularity)}
\\
&=
\frac{\partial\theta}{\partial\theta}
=1.
\end{aligned}
\]
Plug these expectations into the Cauchy--Schwarz inequality and rearrange.
\hfill$\square$

\paragraph{Remark.}
Sometimes, the MLE is unbiased and attains the Cram\'er--Rao lower
bound, but not always. However, we will show this is true asymptotically.

If $\Theta\subseteq\mathbb{R}^p$ with $p>1$, let
\[
\Phi:\Theta\to\mathbb{R}
\]
be a scalar function of $\theta$.
Consider an unbiased estimator $\tilde{\Phi}(X_1,\ldots,X_n)$.
Then one can show that
\[
\operatorname{Var}_{\theta}(\tilde{\Phi})
\geq
\frac{1}{n}
\left(
\frac{\partial\Phi(\theta)}{\partial\theta}
\right)^{T}
I(\theta)^{-1}
\frac{\partial\Phi(\theta)}{\partial\theta},
\qquad
\forall\theta\in\operatorname{int}(\Theta).
\]

\paragraph{Corollary 1.}
If $\tilde{\theta}$ is not unbiased, we have
\[
\operatorname{Var}_{\theta}(\tilde{\theta})
\geq
\frac{
\displaystyle
\frac{\partial}{\partial\theta}
\mathbb{E}_{\theta}\tilde{\theta}
}{
nI(\theta)
},
\qquad
\forall\theta\in\operatorname{int}(\Theta).
\]

\paragraph{Example.}
Take $\Phi(\theta)=\alpha^{T}\theta$ for some
$\alpha\in\mathbb{R}^p$.
In this case,
\[
\frac{\partial\Phi(\theta)}{\partial\theta}=\alpha,
\]
and
\[
\operatorname{Var}_{\theta}(\tilde{\Phi})
\geq
\frac{1}{n}\alpha^{T}I(\theta)^{-1}\alpha.
\]

\subsection*{Example Sheet 1}

We observe a single pair
\[
(X_1,X_2)
\sim
N\left(
\begin{pmatrix}
\theta_1\\
\theta_2
\end{pmatrix},
\Sigma
\right),
\]
where $\Sigma\succ0$ is known.
We wish to estimate $\theta_1$.

\paragraph{Case 1.}
Assume that $\theta_2$ is known.
In this case, the model has one parameter and F.I.\ matrix
\[
I_1(\theta_1).
\]

\paragraph{Case 2.}
Assume that $\theta_2$ is unknown.
For any estimator
\[
\tilde{\theta}
=
\begin{pmatrix}
\tilde{\theta}_1\\
\tilde{\theta}_2
\end{pmatrix},
\]
we have
\[
\tilde{\theta}_1
=
\begin{pmatrix}
1\\
0
\end{pmatrix}^{T}
\tilde{\theta}.
\]
So, if $\tilde{\theta}$ is unbiased, we have the Cram\'er--Rao lower bound
\[
\operatorname{Var}_{\theta}(\tilde{\theta}_1)
\geq
\begin{pmatrix}
1\\
0
\end{pmatrix}^{T}
I(\theta)^{-1}
\begin{pmatrix}
1\\
0
\end{pmatrix}
=
\bigl[I(\theta)^{-1}\bigr]_{11}.
\]

\paragraph{Goal.}
Show
\[
I_1(\theta_1)^{-1}
<
\bigl[I(\theta)^{-1}\bigr]_{11},
\]
unless $\Sigma$ is diagonal ($X_1,X_2$ independent).

\section*{II. Asymptotic Theory of the MLE}

All reasonable estimators are asymptotically unbiased:
\[
\mathbb{E}_{\theta}\tilde{\theta}
\xrightarrow{n\to\infty}
\theta,
\qquad
\forall\theta\in\Theta.
\]

We wish to prove that estimators are \textbf{consistent}:
\[
\tilde{\theta}\xrightarrow{?}\theta
\qquad
\text{as }n\to\infty
\quad
\text{under }P_{\theta},
\qquad
\forall\theta\in\Theta.
\]

For consistent estimators, a reasonable notion of optimality is
\textbf{asymptotic efficiency}:
\[
n\operatorname{Var}_{\theta}(\tilde{\theta})
\xrightarrow{n\to\infty}
I(\theta)^{-1},
\qquad
\forall\theta\in\operatorname{int}(\Theta).
\]

We will show that the MLE is asymptotically efficient and, furthermore,
\[
\hat{\theta}_{\mathrm{MLE}}
\approx
N\left(\theta,n^{-1}I(\theta)^{-1}\right),
\qquad
\forall\theta.
\]
This distributional convergence can be used for statistical inference.

\subsection*{Stochastic Convergence}

\paragraph{Definition 1.}
Let $(X_n:n\in\mathbb{N})$ be a collection of random variables in some
probability space $(\Omega,\mathcal{F},\mathbb{P})$.

\begin{itemize}
    \item We say that $X_n$ converges almost surely to a random variable
    $X$ on $(\Omega,\mathcal{F},\mathbb{P})$, or
    \[
    X_n\xrightarrow{\mathrm{a.s.}}X,
    \]
    if and only if
    \[
    \mathbb{P}\left(
    \left\{
    \omega:
    \left\|X_n(\omega)-X(\omega)\right\|_2
    \xrightarrow{n\to\infty}0
    \right\}
    \right)
    =1.
    \]

    \item We say that $X_n$ converges to $X$ in probability, or
    \[
    X_n\xrightarrow{P}X,
    \]
    if, for every $\varepsilon>0$,
    \[
    \mathbb{P}\left(
    \|X_n-X\|_2>\varepsilon
    \right)
    \xrightarrow{n\to\infty}0.
    \]
\end{itemize}

\begin{center}
\textbf{End of Lecture 3}
\end{center}
\end{document}
